% Reviewed translation: PDF pages 126--133 / printed pages 112--119.
% The original scan is authoritative; environment headings remain English.
\makeatletter
\gdef\@chapapp{\chaptername}
\gdef\thechapter{\arabic{chapter}}
\gdef\theHchapter{\arabic{chapter}}
\gdef\Hy@chapapp{chapter}
\makeatother
\ctexset{
  chapter/name = {第,章},
  chapter/number = \arabic{chapter},
  chapter/numbering = true
}
\setcounter{chapter}{1}
\chapter{原始变量下 Stokes 问题的数值解}\label{chap:stokes-primitive-variables}
\counterwithin{equation}{section}
\renewcommand{\thesection}{\arabic{section}}
\renewcommand{\thesubsection}{\thesection.\arabic{subsection}}
\renewcommand{\theequation}{\thesection.\arabic{equation}}
\renewcommand{\thestatement}{\thesection.\arabic{statement}}
\renewcommand{\thefigure}{\arabic{figure}}
\setcounter{figure}{6}

\MathBookSourcePage{126}
\section{一般逼近}\label{sec:primitive-general-approximation}

Chapter~\ref{chap:stokes-foundations} 的 Section~\ref{sec:abstract-variational-problem}
所讨论的抽象问题很容易得到一种直接的逼近方法. 在合理假设下, 该逼近收敛,
且误差与所涉及空间的逼近误差成正比. 将这种方法用于 Stokes 问题时,
可以得到速度和压力的协调逼近, 尽管逼近速度场通常并非严格无散.
本章其余部分所发展的各种有限元方法都以本节材料为基础.
非协调方法也可以纳入这一框架(参见 Zine [85]), 但为了简洁起见,
我们将其完全略去.

\subsection{一个抽象逼近结果}\label{subsec:primitive-abstract-approximation}
\setcounter{equation}{0}

本小节研究 Chapter~\ref{chap:stokes-foundations} 的
Section~\ref{sec:abstract-variational-problem} 中所分析的抽象变分问题的逼近.
我们沿用相同的记号, 并把问题置于完全相同的情形中. 回顾问题 $(Q)$:

给定 $l\in X'$ 和 $\chi\in M'$, 求 $(u,\lambda)\in X\times M$, 使得
\begin{align}
a(u,v)+b(v,\lambda)&=\langle l,v\rangle
&&\forall v\in X,\label{eq:primitive-abstract-q-first}\\
b(u,\mu)&=\langle\chi,\mu\rangle
&&\forall\mu\in M.\label{eq:primitive-abstract-q-second}
\end{align}
这里, $X$ 和 $M$ 是两个实 Hilbert 空间, $a(\cdot,\cdot)$ 和
$b(\cdot,\cdot)$ 分别是定义在 $X\times X$ 和 $X\times M$ 上的两个连续双线性形式.
与形式 $b(\cdot,\cdot)$ 相联系, 定义线性算子 $B$ 和 $B'$:
\[
\langle Bv,\mu\rangle=\langle v,B'\mu\rangle=b(v,\mu)
\qquad\forall v\in X,\quad\forall\mu\in M,
\]
并置
\[
V(\chi)=\{v\in X;\ Bv=\chi\},
\qquad
V=V(0)=\operatorname{Ker}(B).
\]

\MathBookSourcePage{127}
回顾与问题 $(Q)$ 相联系的问题 $(P)$:

给定 $l\in X'$ 和 $\chi\in M'$, 求 $u\in V(\chi)$, 使得
\begin{equation}\label{eq:primitive-abstract-p}
a(u,v)=\langle l,v\rangle
\qquad\forall v\in V.
\end{equation}
我们保留下列两个保证问题 $(Q)$ 和 $(P)$ 等价且具有唯一解的假设
(参见 Theorem~\ref{thm:abstract-well-posedness} 及其
Corollary~\ref{cor:abstract-v-elliptic-well-posedness}):

存在常数 $\alpha>0$, 使得
\begin{equation}\label{eq:primitive-continuous-ellipticity}
a(v,v)\geq\alpha\|v\|_X^2
\qquad\forall v\in V;
\end{equation}
存在常数 $\beta>0$, 使得
\begin{equation}\label{eq:primitive-continuous-inf-sup}
\inf_{\mu\in M}\sup_{v\in X}
\frac{b(v,\mu)}{\|v\|_X\|\mu\|_M}
\geq\beta.
\end{equation}

令 $h$ 表示趋于零的离散化参数. 对每个 $h$, 令 $X_h$ 和 $M_h$
为满足
\[
X_h\subset X,
\qquad
M_h\subset M
\]
的两个有限维空间. 记 $X_h'$ 和 $M_h'$ 为它们的对偶空间, 对偶范数为
\[
\|l_h\|_{X_h'}=
\sup_{v_h\in X_h}\frac{\langle l_h,v_h\rangle}{\|v_h\|_X},
\qquad
\|\chi_h\|_{M_h'}=
\sup_{\mu_h\in M_h}\frac{\langle\chi_h,\mu_h\rangle}{\|\mu_h\|_M}.
\]
显然,
\[
\|l\|_{X_h'}\leq\|l\|_{X'},
\qquad
\|\chi\|_{M_h'}\leq\|\chi\|_{M'}
\qquad\forall l\in X',\quad\forall\chi\in M'.
\]

与连续情形一样, 对 $a(\cdot,\cdot)$ 和 $b(\cdot,\cdot)$ 定义算子
\[
A_h\in\mathcal L(X;X_h'),
\qquad
B_h\in\mathcal L(X;M_h'),
\qquad
B_h'\in\mathcal L(M;X_h')
\]
如下:
\begin{align*}
\langle A_hu,v_h\rangle&=a(u,v_h)
&&\forall v_h\in X_h,\quad\forall u\in X,\\
\langle B_hv,\mu_h\rangle&=b(v,\mu_h)
&&\forall\mu_h\in M_h,\quad\forall v\in X,\\
\langle v_h,B_h'\mu\rangle&=b(v_h,\mu)
&&\forall v_h\in X_h,\quad\forall\mu\in M.
\end{align*}
严格地说, $B_h'$ 不是 $B_h$ 的对偶算子; 但把 $B_h$ 限制在 $X_h$ 上,
把 $B_h'$ 限制在 $M_h$ 上时, 它们确实互为对偶算子. 此外, 显然有
\[
\|B_hv\|_{M_h'}\leq\|Bv\|_{M'}
\qquad\forall v\in X,
\]
而 $\|A_hv\|_{X_h'}$ 和 $\|B_h'\mu\|_{X_h'}$ 也满足类似不等式.

对每个 $\chi\in M'$, 定义 $V(\chi)$ 的有限维对应物
\[
V_h(\chi)=\{v_h\in X_h;\ b(v_h,\mu_h)=\langle\chi,\mu_h\rangle
\quad\forall\mu_h\in M_h\},
\]
并置
\[
V_h=V_h(0)=\operatorname{Ker}(B_h)\cap X_h,
\]
即
\begin{equation}\label{eq:primitive-discrete-kernel}
V_h=\{v_h\in X_h;\ b(v_h,\mu_h)=0
\quad\forall\mu_h\in M_h\}.
\end{equation}

\MathBookSourcePage{128}
立即可以看到, 通常 $V_h\not\subset V$ 且
$V_h(\chi)\not\subset V(\chi)$, 因为 $M_h$ 是 $M$ 的真子空间.

现在用问题 $(Q_h)$ 逼近问题 $(Q)$:

求 $(u_h,\lambda_h)\in X_h\times M_h$, 使得
\begin{align}
a(u_h,v_h)+b(v_h,\lambda_h)&=\langle l,v_h\rangle
&&\forall v_h\in X_h,\label{eq:primitive-discrete-q-first}\\
b(u_h,\mu_h)&=\langle\chi,\mu_h\rangle
&&\forall\mu_h\in M_h.\label{eq:primitive-discrete-q-second}
\end{align}
与 $(Q_h)$ 相联系, 再定义问题 $(P_h)$:

求 $u_h\in V_h(\chi)$, 使得
\begin{equation}\label{eq:primitive-discrete-p}
a(u_h,v_h)=\langle l,v_h\rangle
\qquad\forall v_h\in V_h.
\end{equation}
由于 $V_h\not\subset V$, 问题 $(P_h)$ 可以看作问题 $(P)$ 的一个外逼近.
这里同样有: 问题 $(Q_h)$ 的任一解 $(u_h,\lambda_h)$ 的第一分量 $u_h$
也是问题 $(P_h)$ 的解. 其逆命题将在下面的 Theorem 中证明.

\setcounter{statement}{0}
\begin{Theorem}\label{thm:primitive-abstract-discrete-error}
$1^\circ)$ 假设下列条件成立:

(i) $V_h(\chi)$ 非空;

(ii) 存在常数 $\alpha^*>0$, 使得
\begin{equation}\label{eq:primitive-discrete-ellipticity}
a(v_h,v_h)\geq\alpha^*\|v_h\|_X^2
\qquad\forall v_h\in V_h.
\end{equation}
则问题 $(P_h)$ 有唯一解 $u_h\in V_h(\chi)$, 并且存在只依赖于
$\alpha^*$, $\|a\|$ 和 $\|b\|$ 的常数 $C_1$, 使得下列``误差界''成立:
\begin{equation}\label{eq:primitive-discrete-velocity-error}
\|u-u_h\|_X
\leq C_1\left\{
\inf_{v_h\in V_h(\chi)}\|u-v_h\|_X
+\inf_{\mu_h\in M_h}\|\lambda-\mu_h\|_M
\right\}.
\end{equation}

$2^\circ)$ 假设 (ii) 成立, 并且还满足:

(iii) 存在常数 $\beta^*>0$, 使得
\begin{equation}\label{eq:primitive-discrete-inf-sup}
\sup_{v_h\in X_h}\frac{b(v_h,\mu_h)}{\|v_h\|_X}
\geq\beta^*\|\mu_h\|_M
\qquad\forall\mu_h\in M_h.
\end{equation}
则 $V_h(\chi)\neq\varnothing$, 并且存在唯一的 $\lambda_h\in M_h$,
使得 $(u_h,\lambda_h)$ 是问题 $(Q_h)$ 的唯一解. 此外, 存在只依赖于
$\alpha^*$, $\beta^*$, $\|a\|$ 和 $\|b\|$ 的常数 $C_2$, 使得
\begin{equation}\label{eq:primitive-discrete-total-error}
\|u-u_h\|_X+\|\lambda-\lambda_h\|_M
\leq C_2\left\{
\inf_{v_h\in X_h}\|u-v_h\|_X
+\inf_{\mu_h\in M_h}\|\lambda-\mu_h\|_M
\right\}.
\end{equation}
\end{Theorem}

\begin{Proof}
$1^\circ)$ 由于 $V_h(\chi)$ 非空, 取 $u_h^0\in V_h(\chi)$, 并求解:

求 $z_h\in V_h$, 使得
\[
a(z_h,v_h)=\langle l,v_h\rangle-a(u_h^0,v_h)
\qquad\forall v_h\in V_h.
\]
由 \eqref{eq:primitive-discrete-ellipticity}, 该问题有唯一解 $z_h$, 因而
\[
u_h=z_h+u_h^0
\]
是问题 $(P_h)$ 的唯一解.

\MathBookSourcePage{129}
现任取 $w_h\in V_h(\chi)$, 则 $v_h=u_h-w_h\in V_h$, 且
\begin{equation}\label{eq:primitive-proof-vh-energy}
a(v_h,v_h)=\langle l,v_h\rangle-a(w_h,v_h).
\end{equation}
由于 $v_h\in X_h$, 可在 \eqref{eq:primitive-abstract-q-first} 中取 $v=v_h$,
再代入 \eqref{eq:primitive-proof-vh-energy}, 得到
\[
a(v_h,v_h)=a(u-w_h,v_h)+b(v_h,\lambda).
\]
又因为 $v_h\in V_h$, 对所有 $\mu_h\in M_h$ 都有 $b(v_h,\mu_h)=0$. 因此
\begin{equation}\label{eq:primitive-proof-vh-error-identity}
a(v_h,v_h)=a(u-w_h,v_h)+b(v_h,\lambda-\mu_h)
\qquad\forall\mu_h\in M_h.
\end{equation}
$a$ 在 $V_h$ 上的椭圆性以及 $a$ 和 $b$ 的连续性给出
\[
\|v_h\|_X
\leq\frac{\|a\|\|u-w_h\|_X+\|b\|\|\lambda-\mu_h\|_M}{\alpha^*}.
\]
所以
\[
\|u-u_h\|_X
\leq\left(1+\frac{\|a\|}{\alpha^*}\right)\|u-w_h\|_X
+\frac{\|b\|}{\alpha^*}\|\lambda-\mu_h\|_M
\]
对所有 $w_h\in V_h(\chi)$ 和 $\mu_h\in M_h$ 成立. 于是
\eqref{eq:primitive-discrete-velocity-error} 成立, 其中
\[
C_1=\max\left(1+\frac{\|a\|}{\alpha^*},\frac{\|b\|}{\alpha^*}\right).
\]

$2^\circ)$ 把 Lemma~\ref{lem:ch4-abstract-div} 用于 $X_h$ 和 $M_h$.
假设 (iii) 表明 $B_h$ 是从 $X_h$ 中的 $V_h^\perp$ 到 $M_h'$ 的同构.
因此 $V_h(\chi)$ 非空, 由第 $1^\circ)$ 部分, 问题 $(P_h)$ 有唯一解 $u_h$.
此外, 由 Corollary~\ref{cor:abstract-v-elliptic-well-posedness},
存在唯一的 $\lambda_h\in M_h$, 使得 $(u_h,\lambda_h)$ 是问题 $(Q_h)$ 的唯一解.

为导出误差界 \eqref{eq:primitive-discrete-total-error}, 先证明
\begin{equation}\label{eq:primitive-discrete-affine-best-approximation}
\inf_{w_h\in V_h(\chi)}\|u-w_h\|_X
\leq\left(1+\frac{\|b\|}{\beta^*}\right)
\inf_{v_h\in X_h}\|u-v_h\|_X.
\end{equation}
任取 $v_h\in X_h$. 与上面一样, 存在唯一的 $z_h\in V_h^\perp$, 使得
\[
B_hz_h=B_h(u-v_h),
\]
且
\[
\|z_h\|_X
\leq\frac{1}{\beta^*}\|B_h(u-v_h)\|_{M_h'}
\leq\frac{\|b\|}{\beta^*}\|u-v_h\|_X.
\]
若置 $w_h=z_h+v_h$, 则
\[
b(w_h,\mu_h)
=b(u-v_h,\mu_h)+b(v_h,\mu_h)
=\langle\chi,\mu_h\rangle
\qquad\forall\mu_h\in M_h,
\]
所以 $w_h\in V_h(\chi)$. 此外,
\[
\|u-w_h\|_X
\leq\|u-v_h\|_X+\|z_h\|_X
\leq\left(1+\frac{\|b\|}{\beta^*}\right)\|u-v_h\|_X.
\]
由于 $v_h$ 任意, 这就证明了 \eqref{eq:primitive-discrete-affine-best-approximation}.

\MathBookSourcePage{130}
还需估计 $\|\lambda-\lambda_h\|_M$. 由
\eqref{eq:primitive-abstract-q-first} 和 \eqref{eq:primitive-discrete-q-first} 得
\[
b(v_h,\lambda_h-\mu_h)
=a(u-u_h,v_h)+b(v_h,\lambda-\mu_h)
\quad\forall v_h\in X_h,\quad\forall\mu_h\in M_h.
\]
因此, 假设 \eqref{eq:primitive-discrete-inf-sup} 给出
\begin{align*}
\|\lambda_h-\mu_h\|_M
&\leq\frac{1}{\beta^*}
\sup_{v_h\in X_h}\frac{a(u-u_h,v_h)+b(v_h,\lambda-\mu_h)}{\|v_h\|_X}\\
&\leq\frac{1}{\beta^*}
\left\{\|a\|\|u-u_h\|_X+\|b\|\|\lambda-\mu_h\|_M\right\}.
\end{align*}
所以
\begin{equation}\label{eq:primitive-pressure-error}
\|\lambda-\lambda_h\|_M
\leq\frac{1}{\beta^*}\left\{
\|a\|\|u-u_h\|_X
+(\beta^*+\|b\|)\inf_{\mu_h\in M_h}\|\lambda-\mu_h\|_M
\right\}.
\end{equation}
于是, 误差界 \eqref{eq:primitive-discrete-total-error} 立即由
\eqref{eq:primitive-discrete-velocity-error},
\eqref{eq:primitive-discrete-affine-best-approximation} 和
\eqref{eq:primitive-pressure-error} 得到.
\end{Proof}

\setcounter{statement}{0}
\begin{Remark}\label{rem:primitive-improved-velocity-error}
不使用 inf-sup 条件 \eqref{eq:primitive-discrete-inf-sup},
也可以略微改进误差界 \eqref{eq:primitive-discrete-velocity-error}.
事实上, 对 \eqref{eq:primitive-proof-vh-error-identity} 应用
\eqref{eq:primitive-discrete-ellipticity}, 得到
\[
\|v_h\|_X
\leq\frac{1}{\alpha^*}\left(
\|a\|\|u-w_h\|_X
+\sup_{v_h\in V_h}\frac{b(v_h,\lambda-\mu_h)}{\|v_h\|_X}
\right).
\]
因此
\begin{equation}\label{eq:primitive-improved-velocity-error}
\begin{split}
\|u-u_h\|_X
\leq{}&\left(1+\frac{\|a\|}{\alpha^*}\right)
\inf_{w_h\in V_h(\chi)}\|u-w_h\|_X\\
&+\frac{1}{\alpha^*}\inf_{\mu_h\in M_h}
\sup_{v_h\in V_h}\frac{b(v_h,\lambda-\mu_h)}{\|v_h\|_X}.
\end{split}
\end{equation}
注意, 表达式
\[
\inf_{\mu_h\in M_h}\sup_{v_h\in V_h}
\frac{b(v_h,\lambda-\mu_h)}{\|v_h\|_X}
\]
反映了 $V_h\not\subset V$ 这一事实: 当 $V_h\subset V$ 时, 它为零.
\end{Remark}

\setcounter{statement}{1}
\begin{Remark}\label{rem:primitive-discrete-optimization}
除假设 \eqref{eq:primitive-discrete-ellipticity} 和
\eqref{eq:primitive-discrete-inf-sup} 外, 再假设双线性形式
$a(\cdot,\cdot)$ 在 $X_h$ 上对称且半正定, 则可以把问题 $(P_h)$ 和 $(Q_h)$
与最优化问题联系起来. 与连续情形一样, 并采用相同记号, 可以证明问题 $(P_h)$
的解 $u_h$ 由下式刻画:
\[
J(u_h)=\inf_{v_h\in V_h(\chi)}J(v_h),
\]
而问题 $(Q_h)$ 的解 $(u_h,p_h)$ 由下式刻画:
\[
\mathcal L(u_h,p_h)
=\min_{v_h\in X_h}\sup_{q_h\in M_h}\mathcal L(v_h,q_h)
=\max_{q_h\in M_h}\inf_{v_h\in X_h}\mathcal L(v_h,q_h).
\]
\end{Remark}

\MathBookSourcePage{131}
\setcounter{statement}{2}
\begin{Remark}\label{rem:primitive-discrete-solution-bound}
由 Theorem~\ref{thm:primitive-abstract-discrete-error} 的论证立即得到:
若假设 \eqref{eq:primitive-discrete-ellipticity} 和
\eqref{eq:primitive-discrete-inf-sup} 成立, 则解 $(u_h,\lambda_h)$ 满足
\[
\|u_h\|_X
\leq\frac{1}{\alpha^*}\left\{
\|l\|_{X'}+\frac{1}{\beta^*}(\alpha^*+\|a\|)\|\chi\|_{M'}
\right\},
\]
\[
\|\lambda_h\|_M
\leq\frac{1}{\beta^*}\left\{\|l\|_{X'}+\|a\|\|u_h\|_X\right\}.
\]
\end{Remark}

只要对所有非零 $v_h$ 都有 $a(v_h,v_h)>0$, 双线性形式
$a(\cdot,\cdot)$ 就在 $V_h$ 上椭圆. 类似地, 只要
$\operatorname{Ker}(B_h')\cap M_h=\{0\}$, 双线性形式 $b(\cdot,\cdot)$
就满足离散 inf-sup 条件 \eqref{eq:primitive-discrete-inf-sup}.
当然, 在这两种情形中, 常数 $\alpha^*$ 和 $\beta^*$ 通常都依赖于 $h$.
而为了在 Theorem~\ref{thm:primitive-abstract-discrete-error} 中得到最优误差界,
显然必须使这两个常数都与 $h$ 无关. 由于通常 $V_h\not\subset V$,
$a(\cdot,\cdot)$ 在 $V$ 上的椭圆性未必能传递到 $V_h$ 上.
因此, 必须在每个具体情形中检验假设
\eqref{eq:primitive-discrete-ellipticity}; 但对我们所考虑的应用而言,
这并不是主要障碍. 另一方面, 离散 inf-sup 条件
\eqref{eq:primitive-discrete-inf-sup} 是 $X_h$ 和 $M_h$ 之间的一致相容性条件,
检验它要困难得多. 下面由 Fortin [28] 给出的 Lemma 为该条件建立了一个实用判据.

\setcounter{statement}{0}
\begin{Lemma}\label{lem:primitive-fortin-criterion}
离散 inf-sup 条件 \eqref{eq:primitive-discrete-inf-sup} 以与 $h$ 无关的常数
$\beta^*>0$ 成立, 当且仅当存在算子 $\Pi_h\in\mathcal L(X;X_h)$, 满足
\begin{equation}\label{eq:primitive-fortin-orthogonality}
b(v-\Pi_hv,\mu_h)=0
\qquad\forall\mu_h\in M_h,\quad\forall v\in X,
\end{equation}
以及
\begin{equation}\label{eq:primitive-fortin-stability}
\|\Pi_hv\|_X\leq C\|v\|_X
\qquad\forall v\in X,
\end{equation}
其中常数 $C>0$ 与 $h$ 无关.
\end{Lemma}

\begin{Proof}
先假设存在这样的算子 $\Pi_h$. 对所有 $\mu_h\in M_h$, 由
\eqref{eq:primitive-fortin-orthogonality} 有
\[
\sup_{v_h\in X_h}\frac{b(v_h,\mu_h)}{\|v_h\|_X}
\geq\sup_{v\in X}\frac{b(\Pi_hv,\mu_h)}{\|\Pi_hv\|_X}
=\sup_{v\in X}\frac{b(v,\mu_h)}{\|\Pi_hv\|_X}.
\]
于是, \eqref{eq:primitive-fortin-stability} 和
\eqref{eq:primitive-continuous-inf-sup} 给出
\[
\sup_{v_h\in X_h}\frac{b(v_h,\mu_h)}{\|v_h\|_X}
\geq\frac{1}{C}\sup_{v\in X}\frac{b(v,\mu_h)}{\|v\|_X}
\geq\frac{\beta}{C}\|\mu_h\|_M.
\]
因此 \eqref{eq:primitive-discrete-inf-sup} 成立, 且 $\beta^*=\beta/C$.

\MathBookSourcePage{132}
反过来, 假设 \eqref{eq:primitive-discrete-inf-sup} 以与 $h$ 无关的常数
$\beta^*>0$ 成立. 对每个 $v\in X$, 存在唯一的
$\Pi_hv\in V_h^\perp$, 使得
\[
b(\Pi_hv,\mu_h)=b(v,\mu_h)
\qquad\forall\mu_h\in M_h,
\]
并且
\[
\|\Pi_hv\|_X
\leq\frac{1}{\beta^*}\|B_hv\|_{M_h'}
\leq\frac{\|b\|}{\beta^*}\|v\|_X.
\]
显然, $\Pi_h\in\mathcal L(X;X_h)$, 并以
$C=\|b\|/\beta^*$ 满足 \eqref{eq:primitive-fortin-stability}.
\end{Proof}

实际中, 构造 $\Pi_h$ 绝非易事. 读者将在 Section~\ref{subsec:primitive-discrete-inf-sup-methods}
中看到如何在若干情形下不显式构造 $\Pi_h$ 而建立 inf-sup 条件.

\setcounter{statement}{3}
\begin{Remark}\label{rem:primitive-inf-sup-riesz-form}
离散 inf-sup 条件 \eqref{eq:primitive-discrete-inf-sup} 还可以方便地写成:
对每个 $\mu_h\in M_h$, 存在 $v_h\in X_h$ (在 $V_h^\perp$ 中唯一), 使得
\[
b(v_h,\mu_h)=\|\mu_h\|_M^2,
\qquad
\|v_h\|_X\leq\frac{1}{\beta^*}\|\mu_h\|_M.
\]
这个结果在连续情形中也成立, 并且明确使用了 $\|\cdot\|_M$ 是 Hilbert 范数这一事实.
\end{Remark}

\setcounter{statement}{4}
\begin{Remark}\label{rem:primitive-pressure-error-inner-product}
在双线性形式 $a(\cdot,\cdot)$ 恰好等于与 Hilbert 范数
$\|\cdot\|_X$ 相联系的标量积 $((\cdot,\cdot))_X$ 的特殊情形中,
公式 \eqref{eq:primitive-pressure-error} 简化为
\begin{equation}\label{eq:primitive-pressure-error-prime}
\|\lambda-\lambda_h\|_M
\leq\frac{1}{\beta^*}\left\{
\inf_{w_h\in V_h(\chi)}\|u-w_h\|_X
+(\beta^*+\|b\|)\inf_{\mu_h\in M_h}\|\lambda-\mu_h\|_M
\right\}.
\tag{1.17$'$}
\end{equation}
事实上,
\[
b(v_h,\lambda_h-\mu_h)
=((u-w_h,v_h))_X+b(v_h,\lambda-\mu_h)
\]
对所有 $v_h\in V_h^\perp$, $w_h\in V_h(\chi)$ 和 $\mu_h\in M_h$ 成立,
而 Remark~\ref{rem:primitive-inf-sup-riesz-form} 中的 $v_h\in V_h^\perp$
给出 \eqref{eq:primitive-pressure-error-prime}.
\end{Remark}

Theorem~\ref{thm:primitive-abstract-discrete-error} 立即给出下列一般收敛结果.

\setcounter{statement}{0}
\begin{Corollary}\label{cor:primitive-general-velocity-convergence}
假设下列条件成立:

$1^\circ)$ 形式 $a(\cdot,\cdot)$ 以与 $h$ 无关的常数 $\alpha^*>0$
满足 \eqref{eq:primitive-discrete-ellipticity};

$2^\circ)$ 存在 $V(\chi)$ 的稠密子集 $\mathcal V(\chi)$,
$M$ 的稠密子空间 $\mathcal M$, 以及两个映射
\[
r_h:\mathcal V(\chi)\longrightarrow V_h(\chi),
\qquad
\rho_h:\mathcal M\longrightarrow M_h,
\]
使得
\[
\lim_{h\to0}\|r_hv-v\|_X=0
\qquad\forall v\in\mathcal V(\chi),
\]
以及
\[
\lim_{h\to0}\|\rho_h\mu-\mu\|_M=0
\qquad\forall\mu\in\mathcal M.
\]
则
\[
\lim_{h\to0}\|u-u_h\|_X=0.
\]
\end{Corollary}

\MathBookSourcePage{133}
\setcounter{statement}{1}
\begin{Corollary}\label{cor:primitive-general-mixed-convergence}
保留上述关于 $a(\cdot,\cdot)$ 和 $M$ 的假设, 并假设
$b(\cdot,\cdot)$ 一致地满足 inf-sup 条件
\eqref{eq:primitive-discrete-inf-sup}. 若存在 $X$ 的稠密子空间
$\mathcal X$ 和映射 $r_h:\mathcal X\to X_h$, 满足
\[
\lim_{h\to0}\|r_hv-v\|_X=0
\qquad\forall v\in\mathcal X,
\]
则
\[
\lim_{h\to0}\left\{\|u-u_h\|_X+\|\lambda-\lambda_h\|_M\right\}=0.
\]
\end{Corollary}
